\begin{tikzpicture}[font=\scriptsize, >=latex]
% dynamic instruction count: mem-bound vs register version (real)
\node[align=left, font=\footnotesize\bfseries] at (4.5,4.6)
  {instructions retired, sum-1..10 (real sim runs)};
\draw[thick, fill=accentsoft] (1.2,1.0) rectangle (3.4,4.0);
\node[font=\footnotesize] at (2.3,2.5) {110};
\node[align=center, font=\tiny] at (2.3,0.7) {memory-bound\\(i, s in RAM)};
\draw[thick, fill=goodbg] (5.6,1.0) rectangle (7.8,3.0);
\node[font=\footnotesize] at (6.7,2.0) {57};
\node[align=center, font=\tiny] at (6.7,0.7) {register version\\(i, s in t0, t1)};
\draw[->, thick, slate] (1.2,1.0) -- (8.6,1.0);
\node[align=left, text=slate, text width=10cm] at (4.8,-0.4)
  {Predicted $\sim$2$\times$, measured 1.93$\times$. The loop body is
   7 instructions vs 3 --- keeping variables in registers is not an
   optimisation, it is the difference between code and bookkeeping.
   Counted by \ttfamily ch03\_rv32sim.py\rmfamily, not estimated.};
\end{tikzpicture}
